\documentclass[10pt]{article}

\usepackage[utf8]{inputenc}

\usepackage[T1]{fontenc}

\usepackage{amsmath}

\usepackage{amsfonts}

\usepackage{amssymb}

\usepackage[version=4]{mhchem}

\usepackage{stmaryrd}

\usepackage{graphicx}

\usepackage[export]{adjustbox}

\graphicspath{ {./images/} }



\begin{document}

Основные свойства Т. ф.: область определения, множество значений, четность и участки монотонности приведены в табл.



Каждая Т. ф. в каждой точке своей области определения непрерывна и бесконечно диффференцируема; пропзводные Т. ф.:



\[

\begin{aligned}

(\sin x)^{\prime} & =\cos x, & (\cos x)^{\prime} & =-\sin x, \\

(\operatorname{tg} x)^{\prime} & =\frac{1}{\cos ^{2} x}, & (\operatorname{ctg} x)^{\prime} & =-\frac{1}{\sin ^{2} x} .

\end{aligned}

\]



Интегралы от Т. ф.:



\[

\int \sin x d x=-\cos x+C, \quad \int \cos x d x=\sin x+C,

\]



$\int \operatorname{tg} x d x=-\ln |\cos x|+C, \quad \int \operatorname{ctg} x d x=\ln |\sin x|+C$.



Все T. ф. допускают разложение в степенные ряды:



$\sin x=x-\frac{x^{3}}{3 !}+\frac{x^{5}}{5 !}-\ldots+(-1)^{n} \frac{x^{2 n+1}}{(2 n+1) !}+\ldots$



при $|x|<\infty$,



$\cos x=1-\frac{x^{2}}{2 !}+\frac{x^{4}}{4 !}-\frac{x^{6}}{6 !}+\ldots+(-1)^{n} \frac{x^{2 n}}{(2 n) !}+\ldots$



при $|x|<\infty$,



$\operatorname{tg} x=x+\frac{1}{3} x^{3}+\frac{2}{15} x^{5}+\frac{17}{315} x^{7}+\ldots$



$\ldots+\frac{2^{2 n}\left(2^{2 n}-1\right)\left|B_{n}\right|}{(2 n) !} x^{2 n-1}+\ldots$ при $|x|<\frac{\pi}{2}$,



$\operatorname{ctg} x=\frac{1}{x}-\left[\frac{x}{3}+\frac{x^{8}}{45}+\frac{2 x^{5}}{945}+\frac{x^{7}}{4725}+\ldots+\frac{2^{2 n}\left|B_{n}\right|}{(2 n) !} x^{2 n-1}\right]$,



при $0<|x|<\pi\left(B_{n}\right.$ - числа Бернулли).



\begin{center}

\includegraphics[max width=\textwidth]{2023_07_09_ec67b0a34343a44edab4g-1}

\end{center}



Функция $y=\sin x$, являющаяся обратной по отношению к функции $x=\sin y$, определяет $y$ как многозначную функцию от $x$; она обозначается $y=\operatorname{Arcsin} x$. Аналогично определяются функции, обратные по отношению к другим Т. ф.; все они наз. обратиыли тригонометрическими бункциями.



Т ритономет рические функциикомІ ле к с ного пе ременн ого. Т. ф. для комплексных значений переменного $z=x+i y$ оиределяются как аналитические продолжения соответствуюпцх Т. ф. действительного переменного в комплексную плоскость.



Так, $\sin z$ и $\cos z$ можно определить с помощью рядов пля $\sin x$ и $\cos x$. Эти ряды сходятся во всей плоскости, поэтому $\sin z$ и $\cos z-$ целые бункции.



Т. ф. тангенс и котангенс определяются формулами



$\operatorname{tg} z=\frac{\sin z}{\cos z}, \quad \operatorname{ctg} z=\frac{\cos z}{\sin z}$.



Т. Ф. $\operatorname{tg} z$ и $\operatorname{ctg} z$-мероморфные функции. Полюсы $\operatorname{tg} z$ простые (1-го порядка) и находятся в точках $z=$ $=\pi / 2+\pi n$, полюсы $\operatorname{ctg} z$ также простые и находятся в точках $z=\pi n, n=0, \pm 1, \ldots$.



Все формулы, справедливые для Т. ф. действитөльного аргумента, остаются справедливыми и для комплексного аргумента.



В отличие от Т. ф. действительного переменного, функции $\sin z$ и $\cos z$ принимают все комплексные значения: уравнения $\sin z=a$ и $\cos z=a$ имеют решения для любого комплексного $a$ :



\[

\begin{aligned}

& z=\arcsin a=-i \ln \left(i a+\sqrt{1-a^{2}}\right), \\

& z=\arccos a=-i \ln \left(a+\sqrt{a^{2}-1}\right) .

\end{aligned}

\]



Т. ф. $\operatorname{tg} z$ и $\operatorname{ctg} z$ принимают все комплексные значения, кроме $\pm i:$ уравнения $\operatorname{tg} z=a, \operatorname{ctg} z=a$ имеют решения для любого комплексного числа $a \neq \pm i$ :



\[

\begin{aligned}

& z=\operatorname{arctg} a=\frac{i}{2} \ln \frac{1-i a}{1+i a}, \\

& z=\operatorname{arcctg} a=\frac{i}{2} \ln \frac{i a+1}{i a-1} .

\end{aligned}

\]



Т. Ф. можно выразить через показапельную бункцию:



\[

\sin z=\frac{1}{2 i}\left(e^{i z}-e^{-i z}\right),

\][



\textbackslash cos z=\textbackslash frac\{1\}\{2\}\textbackslash left(e\^{}\{i z\}+e\^{}\{-i z\}\textbackslash right), \textbackslash quad \textbackslash operatorname\{tg\} z=\textbackslash frac\{1\}\{i\} \textbackslash frac\{e\^{}\{i z\}-e-i z\}\{e\^{}\{i z\}+e\^{}\{-i z\}\},

]



п зиперболические бункции:



$\sin z=-i \operatorname{sh} i z, \quad \cos z=\operatorname{ch} i z, \quad \operatorname{tg} z=-i$ th $i z$.



B. И. Битюиков.



ТРИГОНОМЕТРИЧЕСКИЙ ПОЛИНОМ, К о н е ч-



ная тритонометрическая сумма,выражение вида



\[

T(x)=\frac{a_{0}}{2}+\sum_{k=1}^{n} a_{k} \cos k x+b_{k} \sin k x

\]



с действительными коэффициентами $a_{0}, a_{k}, b_{k}$, $k=1, \ldots, \quad n ;$ число $n$ наз. п о р я д к о м 'Г. П. $\left(\left|a_{n}\right|+\left|b_{n}\right|>0\right)$. Т. $\quad$. можкно записать в комплексной форме



$T(x)=\sum_{k=-n}^{n} c_{k} e^{i k x}$, где $2 c_{k}=\left\{\begin{array}{cc}a_{k}-i b_{k}, & k \geqslant 0, \\ a_{-k}+i b_{-k}, & k<0 .\end{array}\right.$



Т. п. являются важнейшим средством $n p u$ ближения бункций.



B. И. Битючков.



ТРИГОНОМЕТРИЧЕСКИЙ РЯД - ряд по косинусам и синусам кратных дуг, т. е. ряд вида



\[

\frac{a_{0}}{2}+\sum_{k=1}^{\infty}\left(a_{k} \cos k x+b_{k} \sin k x\right)

\]



или в комплексной форме



\[

\sum_{k=-\infty}^{\infty} c_{k} e^{i k x}

\]



где $a_{k}, b_{k}$ или, соответственно, $c_{k}$ наз. к о ә фф ици и н т а ми Т.



Впервые Т. р. встречаются у Л. Эйлера (L. Euler, 1744). Он получил разложения



\[

\begin{gathered}

\frac{\pi-x}{2}=\sin x+\frac{\sin 2 x}{2}+\frac{\sin 3 x}{3}+\ldots(0<x<2 \pi), \\

\frac{1-r \cos x}{1-2 r \cos x+r^{2}}=1+r \cos x+r^{2} \cos 2 x+\ldots, \\

\frac{r \sin x}{1-2 r \cos x+r^{2}}=r \sin x+r^{2} \sin 2 x+\ldots

\end{gathered}

\]



В сер. 18 в. в связи с исследованиями задачи о свободном колебании струны возник вопрос о возможности





\end{document}